\chapter{组合逻辑电路}

组合逻辑的特点很直接：输出只由当前输入决定，不保存过去。它像一个硬件函数，但这个函数不是一步一步执行的过程，而是一张由门和连线组成的网络。输入变化后，信号沿多条路径同时传播；等延迟过去，输出稳定到当前输入对应的值。

\section{一、从真值表到门网络}

先做一个简单判断器：输入 A、B、C 表示三个传感器，要求至少两个传感器为 1 时输出 1。这个需求不能只凭直觉画线，先列出关键行：

\begin{longtable}{p{0.12\linewidth}p{0.12\linewidth}p{0.12\linewidth}p{0.2\linewidth}p{0.32\linewidth}}
\toprule
A & B & C & Y & 原因 \\
\midrule
0 & 0 & 0 & 0 & 一个条件都没有 \\
1 & 0 & 0 & 0 & 只有一个条件 \\
0 & 1 & 0 & 0 & 只有一个条件 \\
0 & 0 & 1 & 0 & 只有一个条件 \\
1 & 1 & 0 & 1 & A 和 B 同时满足 \\
1 & 0 & 1 & 1 & A 和 C 同时满足 \\
0 & 1 & 1 & 1 & B 和 C 同时满足 \\
1 & 1 & 1 & 1 & 三个都满足 \\
\bottomrule
\end{longtable}

从表中可以直接得到一个门网络：

\begin{lstlisting}
Y = (A AND B) OR (A AND C) OR (B AND C)
\end{lstlisting}

这条表达式里有三条“证据路径”。第一条路径证明 A、B 同时为 1；第二条路径证明 A、C 同时为 1；第三条路径证明 B、C 同时为 1。最后的 OR 不是随便合并，而是在问：三条证据里有没有任意一条成立。这样看，门网络就不只是公式变形，而是把真值表中的每一种输出为 1 的原因接成硬件。

\begin{pdfdiagram}[x=1cm,y=1cm]
  \node[hw control,minimum width=0.8cm] (a) at (0,2.2) {A};
  \node[hw control,minimum width=0.8cm] (b) at (0,1.25) {B};
  \node[hw control,minimum width=0.8cm] (c) at (0,0.3) {C};
  \node[hw compute,minimum width=1.3cm] (ab) at (3.0,2.0) {A AND B};
  \node[hw compute,minimum width=1.3cm] (ac) at (3.0,1.25) {A AND C};
  \node[hw compute,minimum width=1.3cm] (bc) at (3.0,0.5) {B AND C};
  \node[hw compute,minimum width=1.1cm] (or) at (6.1,1.25) {OR};
  \node[hw state,minimum width=0.8cm] (y) at (8.3,1.25) {Y};
  \draw[hw bus] (a.east) -- ++(0.8,0) |- (ab.west);
  \draw[hw bus] (b.east) -- ++(0.95,0) |- (ab.west);
  \draw[hw bus] (a.east) -- ++(0.55,0) |- (ac.west);
  \draw[hw bus] (c.east) -- ++(0.8,0) |- (ac.west);
  \draw[hw bus] (b.east) -- ++(0.55,0) |- (bc.west);
  \draw[hw bus] (c.east) -- ++(0.95,0) |- (bc.west);
  \draw[hw bus] (ab.east) -- (or.west);
  \draw[hw bus] (ac.east) -- (or.west);
  \draw[hw bus] (bc.east) -- (or.west);
  \draw[hw bus] (or.east) -- (y.west);
\end{pdfdiagram}
\diagramnote{“至少两个为 1”的组合逻辑由三条证据路径和一个合并路径构成。}

这个例子展示了组合逻辑的一般做法：先把行为说清楚，再把能让输出为 1 的条件提出来，最后用门电路连接。若以后发现路径太长或门太多，可以再优化，但不能跳过“行为完整”这一步。

\section{二、半加器从两个 bit 相加开始}

两个 bit 相加，输出不止一个 bit。因为 1+1 得到二进制 10，需要一个和位 S 和一个进位 C。

\begin{longtable}{p{0.16\linewidth}p{0.16\linewidth}p{0.22\linewidth}p{0.22\linewidth}}
\toprule
A & B & S & C \\
\midrule
0 & 0 & 0 & 0 \\
0 & 1 & 1 & 0 \\
1 & 0 & 1 & 0 \\
1 & 1 & 0 & 1 \\
\bottomrule
\end{longtable}

观察表格可以看到，S 在 A、B 不同时为 1，所以：

\begin{lstlisting}
S = A XOR B
C = A AND B
\end{lstlisting}

这就是半加器。它已经完成了一位相加，但它还缺一个输入：来自低位的进位。只要要做多位加法，最低位之外的每一位都不能只看 A 和 B，还要看低一位是否进上来。

\section{三、全加器把进位接进来}

全加器输入为 A、B、Cin，输出为 S、Cout。可以先把 A 和 B 用半加器相加，得到临时和 S1 和进位 C1；再把 S1 和 Cin 用第二个半加器相加，得到最终和 S 和第二个进位 C2；最后把两个进位合并。

\begin{lstlisting}
S1 = A XOR B
C1 = A AND B
S  = S1 XOR Cin
C2 = S1 AND Cin
Cout = C1 OR C2
\end{lstlisting}

把全加器一位一位串起来，就得到串行进位加法器。它的逻辑很清楚：第 0 位算出进位，传给第 1 位；第 1 位再把进位传给第 2 位。问题也很清楚：若最低位产生的进位一路传到最高位，最高位必须等前面所有进位稳定后才能稳定。

\begin{pdfdiagram}[x=1cm,y=1cm]
  \node[hw compute,minimum width=1.45cm] (fa0) at (1.1,1.1) {全加器 0};
  \node[hw compute,minimum width=1.45cm] (fa1) at (3.6,1.1) {全加器 1};
  \node[hw compute,minimum width=1.45cm] (fa2) at (6.1,1.1) {全加器 2};
  \node[hw compute,minimum width=1.45cm] (fa3) at (8.6,1.1) {全加器 3};
  \draw[hw bus] (fa0.east) -- node[above,hw label] {C1} (fa1.west);
  \draw[hw bus] (fa1.east) -- node[above,hw label] {C2} (fa2.west);
  \draw[hw bus] (fa2.east) -- node[above,hw label] {C3} (fa3.west);
  \draw[hw bus] (0.0,1.1) -- node[above,hw label] {Cin} (fa0.west);
  \draw[hw bus] (fa3.east) -- node[above,hw label] {Cout} (10.0,1.1);
  \node[hw label,align=center,text=BookMuted] at (5.0,0.28) {最坏情况：进位必须从最低位一路传到最高位，最高位结果最后稳定};
\end{pdfdiagram}
\diagramnote{串行进位加法器逻辑简单，但最长延迟路径沿进位链传播。}

这就是组合逻辑里第一次真正遇到“逻辑正确”和“速度足够”的差别。真值表保证最后结果正确，传播延迟决定结果多久之后才可靠。位宽越大，简单串行进位加法器的最坏路径越长。

\section{四、选择器和译码器负责选路}

多路选择器解决“多个输入选一个输出”。2 选 1 选择器的行为是：

\begin{lstlisting}
if sel = 0, Y = A
if sel = 1, Y = B
\end{lstlisting}

这里的 \code{if} 只是描述行为，电路本身没有顺序执行。它的门级表达式是：

\begin{lstlisting}
Y = (NOT sel AND A) OR (sel AND B)
\end{lstlisting}

选择器的作用很基础。若一条输出线可能来自加法器，也可能来自传感器，也可能来自某个寄存器，那么必须有一个明确的选择信号决定哪一路接到输出。没有选择器，多路来源同时驱动同一条线就会冲突。

译码器做相反方向的事：把一个二进制编号展开成 one-hot 信号。2-bit 输入可以选中 4 条输出中的一条。

\begin{longtable}{p{0.18\linewidth}p{0.48\linewidth}}
\toprule
输入 & 输出 \\
\midrule
00 & Y0=1，Y1/Y2/Y3=0 \\
01 & Y1=1，其他为 0 \\
10 & Y2=1，其他为 0 \\
11 & Y3=1，其他为 0 \\
\bottomrule
\end{longtable}

选择器把多路收成一路，译码器把编号展开成多路选择。后面无论是寄存器阵列、数据通路，还是存储器地址选择，都会反复看到这两个结构。

\section{五、组合逻辑没有记忆，也会有短暂毛刺}

组合逻辑只看当前输入，不保存过去。若输入 A 从 0 变成 1，多个门路径的延迟不同，输出可能先短暂抖一下，再稳定到最终值。这种短暂错误叫毛刺。毛刺不一定会造成问题，但前提是后级不要在毛刺期间采样。

这就自然引出下一章。只靠组合逻辑，硬件能计算当前输入对应的输出，却不能决定“什么时候把结果当真”，也不能把某个结果保存到下一步。要让硬件分步骤工作，必须引入能保存状态的电路。

\begin{classicnote}
组合逻辑的验收标准是“输入稳定后，输出在有限延迟后由当前输入唯一决定”。这句话里有三个关键词：当前输入、有限延迟、唯一决定。

\begin{itemize}
\item 纸上实验：给定表达式 \code{Y=(A AND B) OR (NOT A AND C)}，列真值表，标出最长门级路径，并指出输入变化时可能出现毛刺的路径。
\item 位级证据：半加器的 Sum 是 XOR，Carry 是 AND；全加器不是新魔法，而是把进位输入纳入同一张真值表。
\item 结构证据：mux 解决多源选一路，decoder 解决编号选一项；寄存器阵列、内存和数据通路后面都会复用这两个结构。
\item 检查题：若组合网络输出短暂从 0 跳到 1 又回到 0，而最终值正确，它一定是 bug 吗？答案取决于后级是否在毛刺期间采样。
\end{itemize}

经典教材常用“抽象位函数 + 实际延迟模型”一起解释组合电路。只看位函数会漏掉毛刺和关键路径；只看延迟又会失去逻辑行为。两者合在一起，才是硬件设计对象。
\end{classicnote}
